\chapter{Revenu et bénéfice dans l'industrie}\label{ch:in:revenue-and-profit-across-the-industry}

Plaçons côte à côte toutes les entreprises qui publient leurs comptes~: le revenu par employé varie sur plus d'un facteur vingt. En
dollars de 2025, un groupe boursier a dégagé 0.48~million de dollars par tête en 2023, et une société de gestion systématique londonienne 13.0~millions de dollars par tête sur son exercice clos en janvier 2022.
Au sein d'une même entreprise, il varie presque autant~: le \omterm{def:in:revenue-and-profit-across-the-industry:perhead}{revenu par
tête} d'un teneur de marché coté était 4.5 fois plus élevé en 2020, l'année la plus volatile de la période, qu'en 2023, l'une des plus calmes. Le tableau
que ce chapitre construit est le revenu de l'industrie tel que les documents déposés le font apparaître, et ses lacunes font partie du résultat~: les
entreprises qui ne déposent rien comprennent la plupart des teneurs de marché privés décrits dans les chapitres précédents. Le chapitre construit
le tableau à partir de sources, précise ce que mesure et ce que masque un chiffre par tête, mesure comment le \omterm{def:in:revenue-and-profit-across-the-industry:perhead}{revenu par tête} évolue avec la
volatilité du marché pour chaque type d'entreprise, et désigne celles qui manquent.

\section{Un tableau sourcé de chiffres déposés}

\begin{definition}[Revenu par tête, bénéfice par tête]\label{def:in:revenue-and-profit-across-the-industry:perhead}
Le \emph{revenu par tête}\index{revenu par tête} est le revenu d'une entreprise pour une année, après les coûts qu'elle répercute
directement à des tiers, divisé par son effectif~; \emph{bénéfice par tête}\index{bénéfice par tête} est une ligne de résultat publiée
(résultat d'exploitation, résultat avant impôt ou résultat net) divisée par le même effectif. Les deux sont publiés avec la ligne de revenu
ou de résultat, la base de l'effectif (en fin d'exercice ou en moyenne sur l'exercice), la devise et le niveau de prix
utilisé.
\end{definition}

Chaque élément de la définition est un choix qui modifie le résultat, si bien que le tableau les consigne tous. La ligne de revenu retenue est
le chiffre d'affaires du déposant après les coûts pass-through, car une ligne brute gonfle le chiffre des entreprises qui perçoivent des fonds pour le compte d'autres.
Un groupe boursier qui déclare des revenus totaux de 12\,640~millions de dollars pour 2025 indique aussi
des charges liées aux transactions de 2\,709~millions de dollars~: 412~millions de dollars de frais réglementaires (section~31) qu'il perçoit et reverse,
et 2\,297~millions de dollars de paiements de liquidité en espèces, de routage et de compensation. Son revenu après ces postes, 9\,931~millions de dollars, égale ses charges d'exploitation plus son résultat d'exploitation, c'est ainsi que le tableau le calcule à partir des données balisées.
Le revenu d'un teneur de marché est pris après les commissions de courtage, les frais de bourse et de compensation et
les coûts de financement que le livre~16, chapitre~1, qualifie de liés au volume~; celui d'une banque ou d'un courtier est pris après les charges d'intérêts (revenus nets)~; celui d'un gestionnaire de fonds correspond à ses commissions de gestion. L'effectif est celui que le dépôt indique~: la plupart des déposants américains donnent un décompte à
la fin de l'exercice, les sociétés britanniques la moyenne mensuelle de l'exercice,
une banque un chiffre arrondi au millier, et
un teneur de marché européen des équivalents temps plein.

\begin{dated}{2025-12}{Revenu par tête d'après les dépôts, dernière année}\label{dat:in:revenue-and-profit-across-the-industry:table}
\footnotesize
\begin{tabular}{@{}llrrrr@{}}
\toprule
déposant & type & année & revenu, millions & effectif & M\$ par tête\\
\midrule
Nasdaq & bourse & 2025 & 5\,249.0~\$ & 9\,525 & 0.55\\
Man Group & société de gestion & 2025 & 1\,325.0~\$ & 1\,719 & 0.77\\
Intercontinental Exchange & bourse & 2025 & 9\,931.0~\$ & 12\,844 & 0.77\\
Morgan Stanley & banque (groupe) & 2025 & 70\,645.0~\$ & 83\,000 & 0.85\\
Flow Traders & teneur de marché & 2025 & 485.8~\euro & 635 & 0.86\\
MarketAxess & plateforme & 2024 & 817.1~\$ & 891 & 0.94\\
Goldman Sachs & banque (groupe) & 2025 & 58\,283.0~\$ & 47\,400 & 1.23\\
Tradeweb & plateforme & 2025 & 2\,052.4~\$ & 1\,569 & 1.31\\
Coinbase & bourse crypto & 2025 & 7\,181.3~\$ & 4\,951 & 1.45\\
Cboe Global Markets & bourse & 2025 & 2\,429.1~\$ & 1\,661 & 1.46\\
CME Group & bourse & 2025 & 6\,520.6~\$ & 3\,875 & 1.68\\
Interactive Brokers & courtier & 2024 & 5\,185.0~\$ & 2\,998 & 1.77\\
Virtu Financial & teneur de marché & 2025 & 2\,214.9~\$ & 1\,027 & 2.16\\
Optiver & teneur de marché & 2025 & 4\,556.0~\euro & $>2\,000$ & $<2.57$\\
Quadrature Capital & société de gestion & FY25 & 1\,224.8~\pounds & 173 & 9.29\\
\bottomrule
\end{tabular}

\smallskip\noindent Revenu après les coûts pass-through~; effectif en fin d'exercice, sauf Quadrature (moyenne) et
Optiver (borne inférieure indiquée)~; FY25 désigne l'exercice clos en janvier 2025. Par tête, en dollars américains de 2025~: taux de change annuels moyens de la BCE et indice américain des prix à la consommation.
Chaque cellule renvoie à son dépôt dans \texttt{data/industry/filings\_perhead.csv} et \texttt{headcounts.csv}.
\end{dated}

Les quinze déposants forment le panel de ce chapitre, de 2019 à 2025 lorsque les dépôts le permettent~: six bourses et
plateformes de négociation, trois teneurs de marché, deux banques en tant que groupes, deux sociétés de gestion, un courtier et une bourse crypto, soit 87
observations entreprise-année au total, une fois écartées les années sans effectif. Neuf sont lues dans les company facts de la SEC avec
\texttt{firm.filings} (chapitre~11), deux à partir des tableaux dérivés du livre~16, et le reste à partir des rapports annuels, des communiqués de résultats
et du registre britannique, retranscrits avec leurs pages.

Le \omterm{def:in:revenue-and-profit-across-the-industry:perhead}{bénéfice par tête} est plus difficile à regrouper, car la ligne de bénéfice naturelle varie selon le type d'entreprise. Quatre exemples, chacun sur sa
propre ligne~: le résultat d'exploitation de Quadrature s'élevait à 3.22~millions de livres par tête sur son exercice clos en janvier 2025 (4.11~millions de dollars)~;
le résultat d'exploitation de CME Group s'élevait à 1.09~million de dollars en 2025~; le résultat avant impôt de Virtu à 1.07~million de dollars en 2025~; le bénéfice net d'Optiver
s'élevait à 0.65~million d'euros en 2024 (0.71~million de dollars), pour ses 2\,100 employés annoncés.

\section{Par tête~: ce que cela mesure et ce que cela masque}

Le \omterm{def:in:revenue-and-profit-across-the-industry:perhead}{revenu par tête} mesure le revenu que les salariés d'une entreprise produisent ensemble, aux prix propres de l'entreprise, sur son propre
exercice. Ce n'est ni une rémunération ni une productivité au sens où un manager pourrait agir dessus. Ce que le chiffre masque se voit le mieux entreprise par
entreprise (la \cref{fig:in:revenue-and-profit-across-the-industry:ranges}).

\begin{omfigure}
\begin{tikzpicture}
\begin{axis}[width=10.5cm,height=9cm,xmode=log,xmin=0.3,xmax=20,ymin=0.4,ymax=16.6,
  ytick=data,yticklabels={{Nasdaq},{Man Group},{ICE},{Morgan Stanley},{Flow Traders},{MarketAxess},{Goldman Sachs},{Coinbase},{Tradeweb},{Interactive Brokers},{Cboe},{CME Group},{Virtu},{Optiver},{HRT Europe (entité britannique)},{Quadrature}},
  yticklabel style={font=\scriptsize},xticklabel style={font=\footnotesize},
  log ticks with fixed point,xtick={0.5,1,2,5,10},xlabel={\small revenu par tête (millions de dollars de 2025)},
  grid=major,grid style={gray!20},table/col sep=comma]
\addplot+[only marks,mark=*,mark size=1.6pt,color=omDef,error bars/.cd,x dir=both,x explicit,
  error bar style={line width=1.2pt,omDef!45}] table[x=last,y=pos,
  x error plus expr=\thisrow{hi}-\thisrow{last},x error minus expr=\thisrow{last}-\thisrow{lo}]
  {figdata/industry/12-revenue-and-profit-across-the-industry/firm_ranges.csv};
\end{axis}
\end{tikzpicture}

\omcaption{\omterm{def:in:revenue-and-profit-across-the-industry:perhead}{Revenu par tête} des déposants du panel, de 2019 à 2025, en dollars de 2025 sur une échelle logarithmique~: la barre couvre chaque
déposant entre son année la plus basse et son année la plus haute, le point marquant la dernière. Les déposants sont classés selon leur année médiane. L'entité britannique est une
filiale d'un groupe plus vaste, présentée pour comparaison et jamais regroupée. Données~: \texttt{data/industry/filings\_perhead.csv},
via \texttt{in\_industry.firm\_ranges}.}\label{fig:in:revenue-and-profit-across-the-industry:ranges}
\end{omfigure}

\begin{itemize}
\item \textbf{Le modèle d'affaires.} Les chiffres des bourses se concentrent entre 0.5~million et 1.7~million de dollars par tête et
bougent à peine~; la plage d'un teneur de marché est large. L'année la plus haute de Flow Traders vaut 4.5 fois sa plus basse, celle de Virtu 2.5 fois~;
celle de CME Group 1.2 fois.
\item \textbf{Acquisitions.} L'effectif de Nasdaq est passé de 6\,377 à 8\,525 en 2023, «~principalement en raison de notre acquisition
d'Adenza~», et son \omterm{def:in:revenue-and-profit-across-the-industry:perhead}{revenu par tête} est tombé à son plus bas niveau, 0.48~million de dollars~; celui de l'Intercontinental Exchange est passé de
8\,911 à 13\,222 la même année. Une acquisition ajoute des salariés d'un seul coup, et du revenu sur les années suivantes.
\item \textbf{Le groupe entier, et non le desk.} Les deux banques entrent comme groupes entiers, avec leurs effectifs de gestion de fortune, de crédit et
de banque d'investissement. Leurs divisions de marché publient un revenu (chapitre~7) mais aucun effectif, si bien que le
chiffre de la division ne peut pas être calculé à partir des dépôts.
\item \textbf{L'entité, et non le groupe.} La filiale britannique de Hudson River Trading comptait en moyenne 96 salariés en 2022 et 21
en 2025, alors que son revenu après coûts de transaction a augmenté~; son \omterm{def:in:revenue-and-profit-across-the-industry:perhead}{revenu par tête} est passé de 1.00~million de dollars à 7.18~millions de dollars.
Elle comptabilise aussi des produits de prestations facturées à d'autres sociétés du groupe. Les pièges du chapitre~11 s'appliquent~: le tableau exclut ces entités
de tous les chiffres regroupés.
\item \textbf{Le décompte.} Un décompte en fin d'exercice, une moyenne et un décompte en équivalents temps plein divergent le plus lorsqu'une
entreprise grandit ou réduit vite ses effectifs~; «~plus de 2\,000~» est une borne inférieure, si bien que le \omterm{def:in:revenue-and-profit-across-the-industry:perhead}{revenu par tête} qu'il donne est une borne supérieure.
\item \textbf{Devise et niveau de prix.} Le revenu en euros d'une entreprise européenne, converti au taux moyen de l'année, varie avec
le taux de change~; et un dollar de 2019 achetait environ un quart de plus qu'un dollar de 2025, ce que supprime le déflateur.
\end{itemize}

\begin{example}[Normalisation d'une case]\label{ex:in:revenue-and-profit-across-the-industry:cell}
Le revenu net de trading de Flow Traders pour 2020 s'élève à 933.4~millions d'euros, avec 554 équivalents temps plein en fin d'année~:
soit 1.68~million d'euros par tête. Au taux de change moyen de 2020 (1.1422~\$ pour un euro), cela fait 1.92~million de dollars, et au niveau de prix de 2025,
avec un indice des prix à la consommation à 135.83 contre 109.20, soit 2.39~million de dollars. Le chiffre 2023 de la même firme, 300.3
millions d'euros pour 646, devient 0.53~million de dollars. Chaque étape est modeste~; ensemble, elles décident si deux firmes en deux monnaies et deux années
peuvent être comparées du tout.
\end{example}

\begin{method}[Construction du tableau sectoriel]\label{met:in:revenue-and-profit-across-the-industry:table}
\begin{enumerate}
\item Pour chaque déclarant et chaque année, prendre la ligne de revenu après les coûts de pass-through, dans la monnaie du déclarant, avec sa
source~; la dériver de faits balisés lorsque le déclarant ne la présente pas comme une ligne unique.
\item Prendre l'effectif que la déclaration indique, avec sa base, et une phrase citée comme preuve.
\item Convertir en dollars au taux moyen de l'année et au niveau de prix de l'année de base.
\item Diviser~; signaler les bornes inférieures et les entités~; ne jamais combler une année manquante.
\end{enumerate}
\end{method}

\section{Régimes de volatilité~: comment les résultats évoluent avec le marché}

La volatilité du marché fixe le volume de trading et la largeur des fourchettes, et donc une bonne part du revenu de
quiconque tire ses revenus du trading. Le chapitre la mesure par la moyenne annuelle des clôtures quotidiennes de l'indice VIX (tableau du
livre~16)~: 15.39 en 2019, 29.25 en 2020, 19.66 en 2021, 25.64 en 2022, 16.85 en 2023, 15.55 en 2024 et 18.93 en 2025.
La \Cref{fig:in:revenue-and-profit-across-the-industry:series} présente cinq déclarants au fil de ces années.

\begin{omfigure}
\begin{tikzpicture}
\begin{axis}[width=10.5cm,height=6cm,ymode=log,xmin=2018.7,xmax=2025.3,ymin=0.4,ymax=4,
  xtick={2019,2020,2021,2022,2023,2024,2025},xticklabel style={/pgf/number format/1000 sep={},font=\footnotesize},
  log ticks with fixed point,ytick={0.5,1,2,3},yticklabel style={font=\footnotesize},
  ylabel={\small millions de \$ de 2025 par tête},grid=major,grid style={gray!20},table/col sep=comma,
  legend style={font=\scriptsize,legend cell align=left,at={(0.5,-0.18)},anchor=north,/tikz/every even column/.append style={column sep=6pt}},
  legend columns=3]
\addplot[omDef,thick,mark=*] table[x=year,y=Virtu]{figdata/industry/12-revenue-and-profit-across-the-industry/series.csv};
\addplot[omThm,thick,mark=square*] table[x=year,y=Flow_Traders]{figdata/industry/12-revenue-and-profit-across-the-industry/series.csv};
\addplot[black!70,thick,mark=triangle*] table[x=year,y=CME_Group]{figdata/industry/12-revenue-and-profit-across-the-industry/series.csv};
\addplot[omExo,thick,mark=diamond*] table[x=year,y=Goldman_Sachs]{figdata/industry/12-revenue-and-profit-across-the-industry/series.csv};
\addplot[gray,thick,dashed,mark=o,mark options={solid}] table[x=year,y=Nasdaq]{figdata/industry/12-revenue-and-profit-across-the-industry/series.csv};
\legend{Virtu (teneur de marché),Flow Traders (teneur de marché),CME Group (bourse),Goldman Sachs (banque),Nasdaq (bourse)}
\end{axis}
\end{tikzpicture}

\omcaption{\omterm{def:in:revenue-and-profit-across-the-industry:perhead}{Revenu par tête} à travers les régimes de volatilité de 2019--2025 (échelle logarithmique). Les deux teneurs de marché atteignent leur maximum en 2020,
année du VIX moyen le plus élevé, puis reculent d'ici 2023~; les bourses et la banque bougent beaucoup moins. Données~: comme la
\cref{fig:in:revenue-and-profit-across-the-industry:ranges}.}\label{fig:in:revenue-and-profit-across-the-industry:series}
\end{omfigure}

Les médianes par catégorie racontent la même histoire, avec davantage de firmes derrière elles (firmes entières seulement, en millions de dollars de 2025 par tête)~:

\begin{center}\footnotesize
\begin{tabular}{@{}lrrrrrrr@{}}
\toprule
catégorie & 2019 & 2020 & 2021 & 2022 & 2023 & 2024 & 2025\\
\midrule
teneurs de marché & 0.90 & 2.70 & 1.62 & 1.24 & 0.96 & 1.77 & 1.51\\
bourses et plateformes & 1.31 & 1.27 & 1.23 & 1.13 & 1.05 & 1.10 & 1.31\\
banques (firmes entières) & 1.20 & 1.37 & 1.28 & 0.90 & 1.08 & 0.99 & 1.04\\
\bottomrule
\end{tabular}
\end{center}

\noindent La médiane des teneurs de marché a triplé entre 2019 et 2020 et a reculé
d'ici 2023~; celle des bourses a varié dans un écart d'un quart, et celle des banques a été la plus élevée en 2020 et 2021.
(Deux ou trois firmes se trouvent derrière chaque case des teneurs de marché, si bien qu'une médiane y est proche du chiffre d'une seule firme.)

Le \omterm{def:in:revenue-and-profit-across-the-industry:perhead}{revenu par tête} des deux teneurs de marché en 2020 valait 2.16 (Virtu) et 4.51 (Flow Traders) fois leur chiffre de 2023,
tandis que le VIX moyen était 1.74 fois plus élevé. Une régression chiffre ce schéma pour toutes les catégories de firmes à la fois.

\begin{method}[Élasticité du revenu par tête à la volatilité]\label{met:in:revenue-and-profit-across-the-industry:fe}
Ajuster, sur l'ensemble des années-firmes des firmes entières dont l'effectif est ponctuel,
\[
\log\left(\frac{R}{N}\right)_{it}=a_i+b_{k(i)}\log \bar V_t+\varepsilon_{it},
\]
où $R/N$ est le \omterm{def:in:revenue-and-profit-across-the-industry:perhead}{revenu par tête} en dollars constants, $a_i$ un effet fixe pour la firme $i$ (son niveau propre), $\bar V_t$ le
VIX moyen de l'année et $b_k$ l'élasticité des firmes de catégorie $k$. Les effets fixes retirent toute différence de niveau
entre firmes, de sorte que chaque $b_k$ est estimé à partir de la façon dont le chiffre propre de chaque firme varie d'une année sur l'autre. Il faut indiquer $b_k$ avec son
erreur type, le nombre d'années-firmes et de firmes, et la part de la variance intra-firme expliquée.
\end{method}

\begin{omfigure}
\begin{tikzpicture}
\begin{axis}[width=11.6cm,height=5.6cm,xmin=0.4,xmax=6.6,ymin=-2.6,ymax=2,xtick=data,
  xticklabels={{teneur de marché},{bourse ou plateforme},{gestionnaire d'actifs},{banque},{courtier},{plateforme crypto}},
  xticklabel style={font=\scriptsize,text width=1.6cm,align=center},yticklabel style={font=\footnotesize},
  ylabel={\small élasticité au VIX moyen},ytick={-2,-1,0,1,2},grid=major,grid style={gray!20},table/col sep=comma,
  extra y ticks={0},extra y tick style={grid style={black!60,thick}}]
\addplot+[only marks,mark=*,mark size=2pt,color=omDef,error bars/.cd,y dir=both,y explicit,
  error bar style={thick,omDef!60}] table[x=pos,y=b,y error=err]
  {figdata/industry/12-revenue-and-profit-across-the-industry/elasticities.csv};
\end{axis}
\end{tikzpicture}

\omcaption{Élasticité du \omterm{def:in:revenue-and-profit-across-the-industry:perhead}{revenu par tête} au VIX moyen annuel, par catégorie de firme, avec des barres d'erreur à deux erreurs types~: 87
années-firmes de 15 firmes entières, 2019--2025, avec effets fixes par firme. Seule l'élasticité des teneurs de marché est nettement
positive. Données~: comme la \cref{fig:in:revenue-and-profit-across-the-industry:ranges}, via
\texttt{in\_industry.fit}.}\label{fig:in:revenue-and-profit-across-the-industry:elasticity}
\end{omfigure}

L'élasticité des teneurs de marché vaut 1.19, avec une erreur type de 0.26~: une année dont le VIX moyen est supérieur de 10~\% s'accompagne
d'un \omterm{def:in:revenue-and-profit-across-the-industry:perhead}{revenu par tête} supérieur d'environ 12~\%, pour une même firme. Pour les bourses et les plateformes, elle vaut 0.13 (0.16), pour les sociétés
de gestion 0.23 (0.27) et pour les banques, en tant que firmes entières, 0.03 (0.32)~: aucune ne se distingue de zéro. Celle du courtier vaut
$-0.39$ (0.37), et celle de la bourse crypto $-1.27$ (0.59), parce que son revenu suit les cours des crypto-actifs, qui
ont le plus chuté en 2022, année de forte volatilité des actions (chapitre~10). Les pentes expliquent 30~\% de la variance intra-firmes
(\cref{fig:in:revenue-and-profit-across-the-industry:elasticity}).

\begin{remark}[Ce qu'est l'élasticité et ce qu'elle n'est pas]\label{rem:in:revenue-and-profit-across-the-industry:elasticity}
Sept années représentent sept observations du marché~; les erreurs types supposent des résidus indépendants, ce que les années
communes ne permettent pas, si bien qu'elles sont plutôt trop petites. L'élasticité décrit les déclarants du panel sur une période
marquée par un seul grand épisode de volatilité, 2020. Elle décrit la manière dont les catégories de firmes se distinguaient, et non une prévision~: la prochaine année volatile pourra être portée par les taux ou le crédit, là où
un autre ensemble de firmes gagne de l'argent.
\end{remark}

Pour un demandeur d'emploi, la régression a une lecture simple. Le \omterm{def:in:revenue-and-profit-across-the-industry:perhead}{revenu par tête} d'un teneur de marché, et avec lui la part variable
de sa rémunération (chapitre~13), monte et descend avec le marché~; celui d'une bourse, d'une plateforme ou d'une banque y est beaucoup moins sensible.
Les deux sociétés de gestion, dont les commissions suivent les actifs et les résultats de leurs fonds, ne montrent pas non plus de lien net.

\section{Qui manque~: biais de survie et biais de publication}

\begin{definition}[Biais de publication]\label{def:in:revenue-and-profit-across-the-industry:disclosure}
Le \emph{biais de publication}\index{biais de publication} est le biais d'une statistique sur des firmes qui provient de ce que les firmes qui
publient leurs données diffèrent systématiquement de celles qui ne les publient pas~: cotées, régulées ou enregistrées là où les comptes sont
publics, plutôt que privées et enregistrées là où ils ne le sont pas.
\end{definition}

Le haut du tableau est fixé par les firmes qui publient, et le haut de l'industrie par celles qui ne publient pas. Les plus grands
teneurs de marché privés du chapitre~2 ne déposent de compte de résultat consolidé nulle part de manière publique~: leurs broker-dealers américains déposent
des bilans, leurs filiales britanniques déposent des comptes d'entité que le chapitre~11 a montré impossibles à lire comme ceux de la firme, et
les articles de presse sur leurs revenus relèvent, selon la règle de ce livre, du contexte et jamais des données. Un lecteur qui veut le
\omterm{def:in:revenue-and-profit-across-the-industry:perhead}{revenu par tête} de l'industrie au sommet n'a donc aucun chiffre déposé à utiliser, et il doit le dire plutôt que de remplir la case.

Le biais de survie (livre~7, chapitre~3) s'y ajoute. Le panel est construit à partir des firmes qui existent et déposent en 2025~; une firme
rachetée, radiée ou fermée pendant la période ne laisse aucune ligne, et une firme fermée parce que son revenu s'est effondré emporte
avec elle ses pires années. Les deux biais poussent dans des directions connues~: le \omterm{def:in:revenue-and-profit-across-the-industry:disclosure}{biais de publication} écarte les firmes privées les plus rentables,
donc le tableau sous-estime le haut~; le biais de survie écarte les échecs, donc les moyennes avantagent les survivantes. Ce que le
tableau mesure de façon fiable, c'est la forme du revenu de l'industrie, type de firme par type de firme, parmi les firmes que l'on peut vérifier.

\section{Tutoriel~: le tableau de l'industrie}\label{tut:in:revenue-and-profit-across-the-industry:tut}

\begin{tutorial}
\textbf{Objectif.} Construire le panel firme-année à partir des sources, le normaliser et estimer les élasticités.
\textbf{État final~:} le tableau daté,
\cref{fig:in:revenue-and-profit-across-the-industry:ranges,fig:in:revenue-and-profit-across-the-industry:series,fig:in:revenue-and-profit-across-the-industry:elasticity}.
\begin{tutsteps}
\item \textbf{Sources.} \texttt{data/industry/headcounts.csv} contient chaque effectif avec son dépôt et une citation~;
\texttt{revenue\_manual.csv} les déposants relevés dans les rapports et le registre. Le script du chapitre
\texttt{in\_panel\_derive.py}, lancé avec un répertoire de travail provisoire, lit les faits de société de la SEC avec \texttt{firm.filings} et
les tableaux du livre~16, et écrit \texttt{filings\_perhead.csv}~; les tests ne lisent que les fichiers versionnés.
\item \textbf{Normaliser.} \texttt{firm.industrypnl.per\_head(rows, fx, cpi, 2025)} convertit aux taux annuels de la BCE du
chapitre~7 et déflate avec \texttt{data/industry/cpi\_usa.csv}.
\item \textbf{Ajuster.} \texttt{fe\_fit(ph, vix)} construit la matrice de plan des indicatrices de firme et des pentes propres à chaque type de firme
(\cref{lst:in:revenue-and-profit-across-the-industry:fe}).
\omcode{code/firm/industrypnl/firm_industrypnl.py}{88}{113}{Effets fixes par firme et une pente par type de firme, par moindres carrés,
avec erreurs types classiques.}\label{lst:in:revenue-and-profit-across-the-industry:fe}
\item \textbf{Lire.} \texttt{in\_industry.firm\_ranges}, \texttt{spread} et \texttt{mm\_ratios} donnent les chiffres des
sections~2 et~3.
\end{tutsteps}
\end{tutorial}

\textbf{Pour aller plus loin.} Remplacer le VIX par une mesure plus proche du métier de chaque type de firme (un indice de volatilité obligataire pour les plateformes
qui négocient des obligations, les cours des cryptoactifs pour la bourse de crypto-actifs) et comparer les ajustements~; regrouper les erreurs types
par année.

\section{Construction~: le panel de l'industrie}\label{bld:in:revenue-and-profit-across-the-industry:industrypnl}

\begin{build}
\bfield{Rôle} Un panel firme-année sourcé de revenu et d'effectifs, normalisé, pour ce chapitre et pour
la comparaison de la rémunération au revenu du chapitre~14.
\bfield{Interface} \texttt{firm.industrypnl}~: \texttt{Row}~; \texttt{load\_panel}, \texttt{load\_fx},
\texttt{load\_cpi}~; \texttt{real\_usd(row, fx, cpi, base)}~; \texttt{per\_head(rows, fx, cpi, base)}~;
\texttt{ranges(ph)}~; \texttt{ratio(ph, firm, y1, y0)}~; \texttt{fe\_fit(ph, x, kinds)}. NumPy~;
s'appuie sur les taux de \texttt{firm.bankmix}, et le panel est construit avec \texttt{firm.filings}.
\bfield{Règles} Les lignes sans effectif sont écartées, jamais complétées~; le revenu reste dans la devise du déposant tant qu'il n'est pas converti~;
la régression n'utilise que des firmes entières avec des effectifs ponctuels~; les erreurs types sont classiques et présentées comme indicatives.
\bfield{Tests d'acceptation} \texttt{code/firm/industrypnl/tests/}~: sur un panel construit exactement à partir de pentes connues, l'ajustement
les retrouve avec un résidu nul~; avec du bruit, les estimations tombent à moins de quatre erreurs types~; la conversion et
la déflation concordent avec le calcul à la main~; une ligne sans effectif n'est pas chargée.
\bfield{Pour aller plus loin} Erreurs types regroupées par année ou par firme~; une version à effets aléatoires~; le bénéfice opérationnel par tête
lorsque les lignes peuvent être alignées.
\end{build}

\begin{omsources}
\item Formulaires 10-K et faits de société de la SEC de CME Group, Intercontinental Exchange, Nasdaq, Cboe Global Markets,
MarketAxess, Tradeweb, Goldman Sachs, Morgan Stanley et Interactive Brokers, 2019--2025.
\item Flow Traders, rapports annuels 2019 et 2021 et communiqués de résultats pour 2023 et 2025.
\item Hudson River Trading Europe Ltd., comptes pour 2023 et 2025 (Companies House).
\item OECD, indices des prix à la consommation~; BCE, taux de référence de l'euro.
\item Tableaux dérivés du livre~16 pour Virtu Financial, Man Group et le VIX~; chapitres~3, 10 et~11 pour Optiver, Coinbase
et Quadrature Capital.
\end{omsources}

\section{Exercices}

\begin{exercise}[$\star$]\label{exo:in:revenue-and-profit-across-the-industry:1}
Un groupe de bourses publie des revenus totaux de 12\,640~millions de dollars, des charges d'exploitation de 5\,002~millions de dollars et un
résultat d'exploitation de 4\,929~millions de dollars. Calculer ses revenus diminués des charges liées aux transactions, et déterminer le montant de ces charges.
\end{exercise}

\begin{exercise}[$\star$]\label{exo:in:revenue-and-profit-across-the-industry:2}
Dans le tableau daté, repérer les déposants dont le \omterm{def:in:revenue-and-profit-across-the-industry:perhead}{revenu par tête} était inférieur à 1~million de dollars lors de leur dernière année, et calculer
le rapport entre le chiffre le plus élevé et le plus bas parmi les firmes entières.
\end{exercise}

\begin{exercise}[$\star$]\label{exo:in:revenue-and-profit-across-the-industry:3}
Calculer le résultat d'exploitation par tête de CME Group en 2025 et le bénéfice net par tête d'Optiver en 2024, en euros.
\end{exercise}

\begin{exercise}[$\star\star$]\label{exo:in:revenue-and-profit-across-the-industry:4}
Le \omterm{def:in:revenue-and-profit-across-the-industry:perhead}{revenu par tête} de Nasdaq est tombé à son plus bas en 2023. Donner, en s'appuyant sur son 10-K, l'explication qui ne tient pas au marché.
\end{exercise}

\begin{exercise}[$\star\star$]\label{exo:in:revenue-and-profit-across-the-industry:5}
Avec une élasticité de 1.19, calculer de combien le \omterm{def:in:revenue-and-profit-across-the-industry:perhead}{revenu par tête} d'un teneur de marché doit varier entre une année où le VIX moyen
vaut 16.85 et une année où il vaut 29.25. Comparer avec les rapports effectifs de Virtu et de Flow Traders.
\end{exercise}

\begin{exercise}[$\star\star$]\label{exo:in:revenue-and-profit-across-the-industry:6}
Expliquer pourquoi la filiale britannique de Hudson River Trading est écartée de la régression, alors que ses chiffres sont sourcés.
\end{exercise}

\begin{exercise}[$\star\star\star$]\label{exo:in:revenue-and-profit-across-the-industry:7}
\emph{Programmation.} Refaire la régression sans la bourse de crypto-actifs ni le broker. Déterminer comment changent les élasticités et les erreurs types des teneurs de marché et des
bourses, et expliquer pourquoi si peu.
\end{exercise}

\begin{exercise}[$\star\star\star$]\label{exo:in:revenue-and-profit-across-the-industry:8}
\emph{Trouver la faille.} «~Les chiffres déposés montrent que le \omterm{def:in:revenue-and-profit-across-the-industry:perhead}{revenu par tête} du secteur culmine à environ 13~millions de dollars, chez un
gestionnaire d'investissement systématique~; les teneurs de marché sont bien en dessous.~»
\end{exercise}

\section{Problème~: La bonne année}

\begin{problem}[{Problème du week-end --- la bonne année}]\label{pb:in:revenue-and-profit-across-the-industry:1}
Deux cours vendeurs~: l'un d'un teneur de marché coté, l'autre d'un groupe de bourses. Un ami affirme que le teneur de marché «~a de bonnes
années~». Dans quelle mesure, et à quelle fréquence~?

\medskip\noindent\textbf{Partie I --- Le tableau.}
\begin{enumerate}
\item Définir le \omterm{def:in:revenue-and-profit-across-the-industry:perhead}{revenu par tête} et le \omterm{def:in:revenue-and-profit-across-the-industry:perhead}{bénéfice par tête}, et énumérer ce qui doit être publié avec eux.
\item Pourquoi utiliser le revenu après les coûts de pass-through~? Donner l'exemple du groupe de bourses.
\item Nommer les bases sur lesquelles les déclarants indiquent l'effectif et expliquer comment chacune biaise le \omterm{def:in:revenue-and-profit-across-the-industry:perhead}{revenu par tête}.
\item Comment rendre comparables les chiffres en euros et en livres, et ceux d'années différentes~?
\item Combien d'entreprises et d'exercices d'entreprise figurent dans le panel, et d'où proviennent-ils~?
\end{enumerate}

\medskip\noindent\textbf{Partie II --- Ce que cache le tableau.}
\begin{enumerate}[resume]
\item Donner le \omterm{def:in:revenue-and-profit-across-the-industry:perhead}{revenu par tête} le plus bas et le plus élevé à l'échelle de l'entreprise, avec l'entreprise et l'année.
\item Donner le rapport entre l'année la plus élevée et l'année la plus basse pour Flow Traders, Virtu et CME Group.
\item Expliquer les chiffres de Nasdaq et de l'Intercontinental Exchange en 2023.
\item Pourquoi la division marchés d'une banque ne peut-elle pas figurer dans le tableau~?
\item Que s'est-il passé pour le \omterm{def:in:revenue-and-profit-across-the-industry:perhead}{revenu par tête} de la filiale britannique entre 2022 et 2025, et pourquoi ne reflète-t-il pas celui de l'entreprise~?
\end{enumerate}

\medskip\noindent\textbf{Partie III --- Les régimes.}
\begin{enumerate}[resume]
\item Donner la moyenne annuelle du VIX pour 2019--2025.
\item Écrire la régression et dire ce que font les effets fixes par entreprise.
\item Donner l'élasticité et l'erreur type de chaque catégorie.
\item Pourquoi l'élasticité de la bourse de crypto-actifs est-elle négative~?
\item Donner les rapports du \omterm{def:in:revenue-and-profit-across-the-industry:perhead}{revenu par tête} des teneurs de marché entre 2020 et 2023.
\end{enumerate}

\medskip\noindent\textbf{Partie IV --- Le verdict.}
\begin{enumerate}[resume]
\item Énoncer le \emph{résultat nommé}~: l'élasticité du \omterm{def:in:revenue-and-profit-across-the-industry:perhead}{revenu par tête} à la moyenne du VIX par type d'entreprise, avec
les erreurs types, et l'étendue du rapport 2020--2023 parmi les teneurs de marché.
\item Pourquoi les erreurs types sont-elles optimistes~?
\item Définir le \omterm{def:in:revenue-and-profit-across-the-industry:disclosure}{biais de publication} et dire dans quel sens il pousse le haut du tableau.
\item Comment le biais de survie fait-il monter les moyennes~?
\item En deux phrases, répondre à l'ami.
\end{enumerate}
\end{problem}

\section{Questions d'entretien}

\begin{interviewq}[$\star$ \iqroles{trader}]\label{iq:in:revenue-and-profit-across-the-industry:1}
Pourquoi un teneur de marché gagne-t-il davantage une année volatile~? Donnez deux mécanismes.
\end{interviewq}

\begin{interviewq}[$\star$ \iqroles{banque, risque}]\label{iq:in:revenue-and-profit-across-the-industry:2}
Une bourse publie des revenus totaux et des revenus nets des frais liés aux transactions. Lequel utilisez-vous pour la comparer à une
banque, et pourquoi~?
\end{interviewq}

\begin{interviewq}[$\star\star$ \iqroles{chercheur}]\label{iq:in:revenue-and-profit-across-the-industry:3}
Vous régressez le logarithme du \omterm{def:in:revenue-and-profit-across-the-industry:perhead}{revenu par tête} sur le logarithme du VIX, avec effets fixes par entreprise, sur sept ans. Que supprime l'effet fixe,
et que ferait un effet fixe annuel sur votre pente~?
\end{interviewq}

\begin{interviewq}[$\star\star$ \iqroles{chercheur, ingénieur ML}]\label{iq:in:revenue-and-profit-across-the-industry:4}
Votre panel compte quinze entreprises et sept années. Comment calculeriez-vous des erreurs types honnêtes~?
\end{interviewq}

\begin{interviewq}[$\star\star$ \iqroles{développeur, ingénieur ML}]\label{iq:in:revenue-and-profit-across-the-industry:5}
Concevez le stockage d'un panel dont chaque cellule doit porter sa source et dont les sources sont retraitées. Quelles sont les clés~?
\end{interviewq}

\begin{interviewq}[$\star\star\star$ \iqroles{chercheur, trader}]\label{iq:in:revenue-and-profit-across-the-industry:6}
Les plus grands teneurs de marché ne publient pas leurs comptes. Comment estimeriez-vous leur \omterm{def:in:revenue-and-profit-across-the-industry:perhead}{revenu par tête}, et que refuseriez-vous
de conclure~?
\end{interviewq}
